\section*{Materials and methods}

\subsection*{System overview}

MI-Workbench is organised in four layers (Fig~\ref{fig1}A): an interface layer (a command-line interface, a REST application programming interface (API) built on FastAPI, and a WebSocket stream for monitoring); an orchestration layer that compiles a workflow graph and drives the agent loop; a provider-agnostic adapter layer that reaches LLM backends through their command-line interfaces; and a persistence layer (SQLite and a structured artifact directory). All coordination logic---graph compilation, step sequencing, critique merging, stopping and verified execution---sits in the orchestration layer, so the control flow of a run can be followed in one place. Every backend is reached through the same three-method interface, \texttt{run()}, \texttt{smoke\_test()} and \texttt{is\_available()}, so the engine does not depend on any vendor SDK and a deterministic mock backend can replace any provider in tests.

\begin{figure}[!h]
\caption{{\bf MI-Workbench architecture and review cycle.}
(A) The four layers. (B) One cycle of the consensus workflow: the executor's analysis code is run in a sandbox, the artifact and execution report are reviewed concurrently by three reviewer lenses, the lenses' critiques are parsed, merged, escalated on agreement and ranked, and the resulting required fixes are returned to the executor until a stopping condition holds.}
\label{fig1}
\end{figure}

A \emph{workspace} is an isolated research context with its own configuration and budgets; a \emph{run} executes one workflow on one task within a workspace. A run proceeds in \emph{iterations}, each of which is one step of the workflow: one executor call, or one review or consensus step (a whole panel counts as one iteration). One cycle of the consensus workflow, an executor step followed by a panel step, therefore comprises two iterations. Each iteration writes its outputs to \texttt{runs/<run\_id>/iter\_NNNN/}: the step's full output, the artifacts parsed from it (for example \texttt{MECH.md} and \texttt{EVAL.md}), for consensus steps each lens's output and a machine-readable merge record, and for executed steps the full execution report. For every iteration the database and \texttt{run\_meta.json} record the provider, model, reasoning-effort setting, CLI version, input, cached-input and output tokens, duration, grade and, for the run, the reason it stopped. A reproducibility package bundles these files with the resolved prompts, the software environment and the git state. Further system details are in S1 Text.

\subsection*{Workflow graphs}

Workflows are directed graphs declared in YAML. Nodes are agent roles with a prompt template; edges carry one of three conditions: \texttt{always}; \texttt{every\_k:N}, which runs the target in every $N$-th cycle of the plan; and \texttt{on\_flag:F}, which runs the target when the run flag \texttt{F} is set. Before execution the graph is validated (node references, condition syntax and stop-condition values) and compiled into an ordered plan that repeats until a stop condition holds. Nodes with the role \texttt{consensus\_merger} absorb their incoming \texttt{always} edges into a single concurrent panel step. The consensus workflow, bundled as the \texttt{reviewer\_consensus} preset and used in the stopping calibration, the case-study agent runs and the control-plane overhead measurements, is (node settings omitted):

\begin{lstlisting}
nodes:
  - {id: executor, role: executor, prompt_ref: executor/mi_executor}
  - {id: reviewer, role: reviewer, prompt_ref: reviewer/mi_reviewer}
  - {id: adversarial_reviewer, role: adversarial_reviewer,
     prompt_ref: adversarial_reviewer/adversarial_reviewer}
  - {id: bio_plausibility_checker, role: bio_plausibility_checker,
     prompt_ref: biological_plausibility/bio_plausibility_checker}
  - {id: consensus_merger, role: consensus_merger}
edges:
  - {source: executor, target: reviewer, condition: always}
  - {source: executor, target: adversarial_reviewer, condition: always}
  - {source: executor, target: bio_plausibility_checker, condition: always}
  - {source: reviewer, target: consensus_merger, condition: always}
  - {source: adversarial_reviewer, target: consensus_merger, condition: always}
  - {source: bio_plausibility_checker, target: consensus_merger, condition: always}
  - {source: consensus_merger, target: executor, condition: always}
\end{lstlisting}

Four presets are bundled: \texttt{executor\_reviewer} (executor and rigour reviewer, with the adversarial reviewer in every third cycle), \texttt{reviewer\_consensus} (the panel workflow above), \texttt{research\_followups} (executor and a follow-up idea generator) and \texttt{example\_custom} (conditional adversarial escalation, knowledge extraction and loop-level stop conditions).

\subsection*{Agent roles and prompts}

Nine roles are available (Table~\ref{table_roles}), each configured through a role-specific prompt template. A review loop needs only the executor and one or more reviewer lenses, and all evaluations in this paper use only these four roles; the five generative and knowledge roles are optional and were not evaluated. The three lenses emphasise different failure modes: the rigour lens checks claim validity, reproducibility, confounders, negative controls and biological plausibility; the adversarial lens probes seven threat models (p-hacking and forking paths, causal over-claiming, causal confusion, leakage and circularity, biological implausibility, statistical theater and inadequate null models); and the biological-plausibility lens checks tissue and cell-type context, pathway consistency, expression plausibility and regulatory direction. The lenses share some items by design (plausibility appears in all three). A combined-checklist reviewer prompt contains every item of the three lens prompts verbatim; it is the strongest single-reviewer baseline in the evaluation.

\begin{table}[!ht]
\begin{adjustwidth}{-2.25in}{0in}
\centering
\caption{{\bf Agent roles.} The evaluations reported here use only the first four roles; the remaining five are optional and were not evaluated.}
\label{table_roles}
\begin{tabular}{@{}lp{10.5cm}@{}}
\hline
\textbf{Role} & \textbf{Function} \\ \hline
Executor & Performs and revises the analysis; with verified execution, writes one analysis script whose outputs are reported \\
Reviewer (rigour lens) & Claim validity, reproducibility, confounders, negative controls, biological plausibility \\
Adversarial reviewer & Seven threat models for common failure modes of interpretability analyses \\
Biological-plausibility checker & Tissue and cell-type context, pathway consistency, expression plausibility, regulatory direction \\
Idea generator & Prioritised follow-up proposals with value, feasibility and novelty scores \\
Knowledge extractor & Claims, uncertainty estimates and falsification tests (S1 Text) \\
Automator & Reusable pipelines from validated protocols \\
Publication manager & Manuscript planning and checklists \\
Article publisher & Summaries for communication \\ \hline
\end{tabular}
\end{adjustwidth}
\end{table}

Prompts are assembled by role. Only executor steps receive revision framing: the reviewers' feedback, the executor's previous submission and, when execution is enabled, its execution report. Every other step receives the task as context, the iteration number, and the executor's latest full output, with its execution report, under an ``artifact under review'' header.

\subsection*{Consensus review}

A consensus step dispatches the panel concurrently over the current artifact and combines the lenses' critiques in four stages (Fig~\ref{fig1}B).

\textbf{Parsing.} Every reviewer prompt ends with a required machine-readable block: a fenced JSON object listing critiques, each with a \texttt{severity} (critical, high, medium, low or info), a \texttt{category}, a \texttt{description} and a \texttt{required\_fix}. The parser reads this block first, tolerating common JSON defects (trailing commas, single quotes, truncation), then falls back to YAML-like and bracketed formats, and counts each critique once. A lens whose call fails, or whose output contains no readable critique block, is marked failed rather than read as a clean review. A failed panel is retried once and otherwise stops the run with a recorded reason. A partial panel (some lenses usable) is retried once for its failed lenses; if it is still partial, it proceeds, but by default its grade does not enter stopping decisions.

\textbf{Merging.} Critiques from different lenses that describe the same problem are grouped. Four methods are implemented: word-level Jaccard similarity; TF-IDF cosine similarity over normalised, stop-word-filtered and stemmed tokens; sentence-embedding cosine similarity~\cite{reimers2019sbert} (optional); and LLM adjudication, in which one additional call partitions the numbered critiques into groups that identify the same underlying problem. For the similarity-based methods, critiques are grouped by average-linkage agglomerative clustering at a threshold $\tau$, which makes the result independent of input order. By default, critiques from the same reviewer call are not merged with each other, for every merge method: every lens prompt asks for one critique per distinct problem, so two items from one call are kept as separate required fixes; this rule is configurable. An adjudicator answer must be a valid partition of the critique indices; otherwise it is retried once, and if it is still invalid, TF-IDF grouping is used. Based on the merge evaluation (Results), LLM adjudication is the default method, and the thresholds of the Jaccard and TF-IDF methods are set to $\tau=0.1$, the upper (more conservative) end of the range selected by cross-validation (0.075--0.1); in-sample F1 differs by less than 0.001 between these two values.

\textbf{Escalation.} Each group takes the highest severity among its members. If at least $k$ distinct lenses contributed to a group (default $k=2$), its severity is raised by one level, capped at critical. A merged critique keeps the text and proposed fix of every member.

\textbf{Ranking and grading.} Groups are ordered by severity, an optional per-lens weight, the number of contributing lenses and the critique text, so the order never depends on dispatch or parsing order. A letter grade summarises the merged list (two or more critical critiques give F; one critical or three or more high give D; one or more high give C; any medium gives B; otherwise A). Because LLM reviewers raise high-severity concerns about essentially every analysis (Results), the grade is reported as an ordinal summary for monitoring; it is not used as an absolute acceptance criterion by default.

\textbf{Gate.} An optional advancement gate keeps the loop revising while an unresolved critical critique remains; budgets and the iteration cap still bound the run, so a panel that never agrees terminates.

\subsection*{Feedback to the executor}

Merged critiques are stratified by severity: critical and high critiques become required fixes, and medium and lower critiques become suggested improvements. Each entry lists the severity, category, any inferred artifact and section reference, the critique, the reviewer's required fix and the points of merged duplicates. The executor is asked to make targeted revisions to the referenced artifacts rather than rewrite them.

\subsection*{Stopping}

A run stops at the first of: user cancellation; the iteration cap; the token or cost budget (input and output tokens and cost as reported by the provider); a loop stop condition (\texttt{grade\_at\_least:G}, \texttt{on\_flag:F}); or the stopping rule. The stopping rule evaluates, after a minimum number of iterations (default 4, i.e.\ two executor--panel cycles of the consensus workflow), three signals over a window of $w$ observations (default 3): the merged grade is unchanged; no critical or high critique was raised; and consecutive executor outputs have word-level Jaccard similarity of at least 0.9. The combination of signals is configurable (any, all, $k$ of $n$, with optional required signals). Failed panels never enter the signal history, and partial panels do not by default. In addition, a revision budget limits the number of executor revisions after the initial submission; when it is reached, the remaining review steps of the cycle run, so the final artifact is reviewed, and the loop stops. Based on the stopping calibration (Results), the default is a budget of four revisions with the adaptive rule disabled; the adaptive rule can be enabled per run. The reason a run stopped is recorded.

\subsection*{Verified execution}

When verified execution is enabled, the executor is asked for one self-contained Python block that reads only the data paths named in the task, prints a summary of its results, writes \texttt{results.json} and reports no number that the code does not compute. After each executor step the block is run, and the execution report---exit status, standard output, the tail of standard error, and the files written, including \texttt{results.json}---is appended to the artifact the reviewers see and returned to the next executor iteration. With execution enabled, the executor's own agent tools are disabled by default, so every reported number has to come from the executed code. Verification here concerns provenance: it ensures that reported numbers were computed by the executed code, not that the code, its numbers or the conclusions drawn from them are correct.

Execution is isolated by one of three backends. All three run \texttt{python -I} in a fresh temporary directory with a minimal environment, standard input closed, output capped and a wall-clock timeout. The \texttt{-I} flag only isolates the interpreter from \texttt{PYTHON*} environment variables and the user's site-packages; it is not a security boundary. The \texttt{subprocess} backend adds CPU-time, file-size and open-file limits and, on Linux, an address-space limit, runs the code in a new session and kills the whole process group on timeout; it does not restrict network or file-system access. The \texttt{sandbox\_exec} backend adds kernel-enforced confinement on macOS through a Seatbelt profile. The profile denies network access, permits writes only in the temporary directory, and limits reads to system and interpreter paths and to explicitly listed read-only data directories. Data paths that are too broad (the file-system root, the home directory or volume roots) are rejected. The \texttt{docker} backend runs the code in a container without network access, with a read-only root file system, memory, CPU and process limits, dropped capabilities and a non-root user. On macOS, whose kernel does not enforce address-space limits, only the \texttt{docker} backend limits memory. If the requested backend is unavailable, execution fails closed. Automated tests verify network denial, write confinement, read allow-listing, CPU limits and process-group termination for the \texttt{sandbox\_exec} backend by running code under the sandbox; tests of the \texttt{docker} backend check only the container command it constructs, and this backend was not used in the experiments reported here.

\subsection*{Adapters}

The Claude Code, Codex and Gemini CLIs are wrapped as adapters, and a mock adapter returns deterministic, role-aware responses for tests. Each adapter forwards the model and, where the CLI supports it, the reasoning-effort setting (the Gemini CLI has no effort flag, so a requested effort is recorded as unsupported), sends the prompt on standard input, runs the CLI in its own process group and parses the CLI's JSON output for the answer, the token counts and, where reported, the cost. Reviewer calls run with the CLI's agent tools disabled (for Codex: a read-only sandbox, no shell, no web search, no plugins or external tool servers, and the user configuration file, \texttt{config.toml}, ignored). Transient errors (rate limits, timeouts, network and overload errors) are retried with exponential backoff; authentication, quota and configuration errors are not. Sampling temperature is not exposed by these CLIs, so runs use each provider's default decoding; the model, effort setting and CLI version of every call are recorded.

\subsection*{Planted-flaw critique benchmark}

To measure what automated reviewers catch, we built a benchmark of analysis write-ups with known flaws (S1 File). Twelve scenarios cover common interpretability analyses of single-cell foundation models: attention-based network recovery in Geneformer and scGPT, linear probing, in-silico perturbation, sparse-autoencoder features, head ablation, activation patching, cross-tissue threshold transfer, age signals in embeddings, attention-derived hub genes, perturbation-response prediction and gene-embedding similarity. Each scenario has a flawed write-up of 750--1,100 words (question, data and model, methods, numerical results, interpretation, limitations) containing four planted flaws, one from each of four of the five families in Table~\ref{table_taxonomy}; each of the twelve flaw types occurs exactly four times (48 flaws). Three of the four attention-mechanics (T1) flaws are the claim that dot-product attention is symmetric. Six scenarios also have a clean counterpart in which all four flaws are corrected and no instance of any flaw type is present. Flaws are stated naturally, are detectable from the text alone and never appear among the stated limitations. For every flaw the ground truth records the verbatim passage, a description and a strict detection criterion. The items and their ground truth were written with Claude Opus 5.5 (Anthropic), used interactively outside the platform (see Use of language models outside the evaluated system). A separate instance of the same model, working in a fresh context, checked them for verbatim quotes, unplanned flaws, factual errors in non-planted statements and internal numerical consistency. Corrections were made before the benchmark was frozen (SHA-256 manifest) and before any reviewer call. The authoring model therefore comes from a different provider than all reviewer, judge and adjudicator models (OpenAI).

\begin{table}[!ht]
\begin{adjustwidth}{-2.25in}{0in}
\centering
\caption{{\bf Flaw taxonomy of the planted-flaw benchmark.} Each type occurs in four flawed items.}
\label{table_taxonomy}
\begin{tabular}{@{}llp{9.5cm}@{}}
\hline
\textbf{Family} & \textbf{Type} & \textbf{Example of the planted error} \\ \hline
Statistics & S1 multiple comparisons & significance claimed from uncorrected per-head $p$-values across 144 heads \\
 & S2 effect-size neglect & a negligible effect with a tiny $p$-value presented as substantive \\
 & S3 metric misuse & AUROC as the sole metric for edges at 0.3\% density \\
Confounding & C1 expression or co-expression & gene-pair scores interpreted as regulation without co-expression or expression baselines \\
 & C2 reference-degree bias & recovery against a hub-dominated reference without a degree-preserving null \\
 & C3 leakage & donors shared between training and test, or selection on the test set \\
Causal logic & L1 causal over-claiming & correlational attention evidence described as the model's mechanism \\
 & L2 circular validation & validation evidence that is not independent of how the finding was obtained \\
Technical & T1 attention mechanics & ``dot-product attention is symmetric, so $A_{ij}=A_{ji}$'' \\
 & T2 model or tokenization & wrong tokenization scheme, training objective or layer count for a named model \\
Biology & B1 context implausibility & a regulator claimed active in a cell type where it is not expressed \\
 & B2 regulatory logic & reversed regulatory direction or a misattributed pathway \\ \hline
\end{tabular}
\begin{flushleft} Type codes are those used in the benchmark files and results (S1 File); they are unrelated to the reviewer conditions C1--C5 and to the numbering of the Supporting Information items. AUROC, area under the receiver operating characteristic curve.
\end{flushleft}
\end{adjustwidth}
\end{table}

\subsection*{Evaluation protocol}

\textbf{Conditions.} For every item, model configuration and repeat we made eight independent reviewer calls on the same input: three with the rigour lens, three with the combined-checklist reviewer, one with the adversarial lens and one with the biological-plausibility lens. Five conditions were built from these calls: C1, one rigour reviewer; C2, one combined-checklist reviewer; C3, three rigour samples; C4, three combined-checklist samples; and C5, the three-lens panel (one rigour, one adversarial and one biological-plausibility call). The conditions share calls: the call of C1 is one of the three calls of C3 and the rigour call of C5, and the call of C2 is one of the three calls of C4. Because a flaw counts as found if any call of a condition identifies it, C3 and C5 cannot have lower recall than C1 on any item, nor C4 than C2; these contrasts measure how much the additional calls add, whereas the call-matched contrasts C5--C3 and C5--C4 test whether distinct lenses add recall. Conditions C3--C5 use three calls each, so the panel is compared with a self-ensemble of a single lens (C3) and with a self-ensemble of the strongest single reviewer (C4) at the same number of calls; token counts are reported per condition. Reviewers received a fixed task statement and the artifact, with tools disabled. Severity outcomes (merged grade, severity score, severity-aware recall and high or critical critiques on clean items) used one merge setting for all conditions, fixed with the analysis plan before any reviewer call: lexical Jaccard merging at $\tau=0.5$. This threshold merges few paraphrased duplicates (see Merging paraphrased critiques), so these outcomes approximate counts over unmerged critiques. The same outcomes with TF-IDF merging at $\tau=0.075$, the threshold selected in 11 of the 12 cross-validation folds of the merge evaluation, are reported as a sensitivity analysis (S2 Table).

\textbf{Models.} We evaluated three model configurations from one provider (OpenAI) through the Codex CLI (version 0.145.0): gpt-5.6-sol at medium reasoning effort, gpt-5.5 at medium effort and gpt-5.6-luna at low effort, each with three repeats (1,296 reviewer calls; S1 Table). Codex prepends to every prompt the built-in harness instructions of the installed CLI version and the user-level instruction file (\texttt{AGENTS.md} in the Codex home directory), which the CLI reads even when \texttt{config.toml} is ignored. The instruction file is a generic persona for agentic software engineering of about 4\,kB: it asks the model to name inconsistencies and unclear specifications instead of guessing, to point out problems rather than agree, to prefer simple code and to avoid changes outside the task. It contains nothing about biology, interpretability, the benchmark items or the review task. Both texts were identical for all calls, conditions and model configurations, including judge, executor and adjudication calls, and their SHA-256 hashes are recorded for every call.

\textbf{Judging.} For each item, configuration and repeat, one judge call received the item text, its planted flaws with their detection criteria, and the union of all parsed critiques from the eight calls, shuffled with a fixed seed and stripped of lens, call and severity. The judge labelled every critique either with the planted flaw whose detection criterion it satisfied or as unmatched; unmatched critiques were further labelled substantive (a specific, plausible concern not in the ground truth), generic (boilerplate) or incorrect (factually wrong, or misreading the analysis). Two judge models were used (gpt-5.6-sol and gpt-5.5, both at high effort); both are also evaluated as reviewers, at medium effort. Agreement is reported as Cohen's $\kappa$~\cite{cohen1960kappa}; the first judge is primary, and all headline results were recomputed with the second judge and with the labels on which both judges agreed.

\textbf{Outcomes and statistics.} The primary outcome is flaw recall: the fraction of an item's planted flaws identified by at least one critique of the condition. Secondary outcomes are recall by family and type, severity-aware recall (the flaw is identified by a merged critique of high or critical severity), the flaws found by only one panel lens, discrimination between flawed and clean items, measured by the AUROC of the merged grade and of a severity score that sums 8, 4, 2, 1 and 0 points for critical to info critiques, high and critical critiques on clean items, the composition of unmatched critiques, and tokens. Five contrasts were planned (C5--C1, C5--C3, C5--C4, C4--C2 and C3--C1) and tested with a paired bootstrap over items (10,000 resamples, repeats kept within item), with Holm correction~\cite{holm1979simple} within each configuration and for the pooled analysis. The analysis plan was fixed before any reviewer call (S1 File); raw reviewer outputs, judge labels, analysis scripts and processed results are in S2 File.

\subsection*{Merge evaluation}

Within each panel instance, two critiques from different lenses that the primary judge matched to the same planted flaw form a duplicate pair, and two critiques matched to different flaws form a non-duplicate pair; unmatched critiques are excluded. Each merging method was scored by pairwise precision, recall and F1 score (their harmonic mean), with its threshold selected by leave-one-scenario-out cross-validation (maximising F1 on the other scenarios). The method with the best held-out F1 was pre-specified as the platform default. We also measured escalation validity: the proportion of multi-lens groups that correspond to a planted flaw. Adjudications, analysis code and results are in S2 File.

\subsection*{Stopping calibration}

To calibrate the stopping rule, we ran real executor--panel loops on the twelve flawed items for a fixed horizon without early stopping: the panel reviewed the original write-up, the executor revised it, and so on for five revisions, in two configurations (gpt-5.6-sol at medium effort and gpt-5.6-luna at low effort, for both executor and panel). Panels merged critiques with lexical Jaccard similarity at $\tau=0.5$, the same setting as the benchmark's severity outcomes; at this threshold few cross-lens duplicates are merged (Results), so the executor received the lenses' critiques largely unmerged and few agreements were escalated. Because no data were available, the executor could correct or withdraw claims, add caveats and specify analyses required before a claim could stand, but was instructed not to report new numerical results. After every revision a single judge (gpt-5.6-sol at high effort; the same model as the executor and panel of the first configuration) labelled each planted flaw as resolved, unresolved, or resolved by fabrication (the revision asserts results or analyses that were not performed) and listed newly introduced errors. Each candidate stopping rule was then replayed on the recorded trajectories with the platform's own replay function, which applies the rule exactly as the engine does, giving its stopping point, whether it stopped while a planted flaw was unresolved (a premature stop) and the number of revisions used. The candidates were each of the three signals alone, any of them, all of them, and the absence of critical or high critiques combined with either grade stability or output similarity (with and without the advancement gate), all with the default window and minimum number of iterations; variants with a window of two and with similarity thresholds of 0.8 and 0.95; and fixed budgets of one to five revisions. The pre-specified selection rule chooses the rule with the lowest premature-stop rate among rules that stop before the horizon in at least half of the runs, breaking ties by fewer revisions. The analysis plan does not state whether the rule is applied within each configuration or to all runs together; because the platform has a single default, we applied it to the 24 runs pooled over both configurations and also report the selection within each configuration. The criterion counts only planted flaws: errors introduced by the revisions are reported for every rule but do not enter the selection. Loop trajectories, state judgments and replay code are in S2 File.

\subsection*{Executed case study}

\textbf{Question and data.} We asked whether Geneformer attention encodes transcription-factor (TF)$\to$target regulation beyond co-expression, expression-rank proximity and hub (degree) structure. One thousand cells were sampled from the immune compartment of Tabula Sapiens~\cite{tabulasapiens2022} (10x 3$'$ v3 chemistry only, stratified by cell type, fixed seed) and tokenised from raw unique molecular identifier (UMI) counts with the Geneformer V2 rank-value convention (normalisation by total counts and by each gene's non-zero median in the pretraining corpus, descending rank, at most 4,096 tokens). Geneformer V2-104M (12 layers, 12 heads; Hugging Face snapshot fcd26c45)~\cite{theodoris2023transfer} was run in evaluation mode, and for a set of 1,322 genes (the 1,200 most frequently tokenised genes plus the TRRUST TFs present in at least 30\% of cells) the attention from each query gene to each key gene was averaged over the cells in which both genes were present, per layer and per head. The resulting data kit also contains Pearson and Spearman co-expression, mean rank distances, co-presence counts and the TRRUST~\cite{han2018trrust} and DoRothEA (confidence A--C)~\cite{garciaalonso2019dorothea} edges among these genes (424 and 2,118 edges, excluding self-loops). Construction scripts, file checksums, package versions and sanity checks are in S3 File.

\textbf{Agent runs.} The \texttt{reviewer\_consensus} workflow was run on the kit with verified execution (\texttt{sandbox\_exec} backend; kit readable, network denied, 900-s wall-clock, 1,800-s CPU and 64-MiB file-size limits) for five executor iterations (an initial analysis and four revisions, the default revision budget), each followed by a panel review, with gpt-5.6-sol (two replicates) and gpt-5.5, all at medium effort and with LLM-adjudicated merging. A control run with gpt-5.6-sol and execution disabled measured what an executor reports without executing code. The task text described every kit file but contained no results. Convergence was disabled, so every run used this full horizon. Executor iterations are numbered 1--5 within each run, and each is followed by its panel review. For each executor iteration we extracted every reported numerical result, checked whether it appears, at the reported precision, in the outputs of that or an earlier execution, and recorded whether it also occurs as a literal in the executed script.

\textbf{Independent analysis.} Separately from the platform's agents, we analysed the same kit with our own code, written with LLM assistance (see Use of language models outside the evaluated system; S3 File). We measured attention symmetry per layer and head (correlation of $A_{ij}$ with $A_{ji}$ and relative asymmetry) and directionality on reference edges (paired comparison of TF-as-query with target-as-query attention). We measured recovery of TRRUST edges among all ordered TF--gene pairs by AUROC and by the area under the precision--recall curve (AUPRC), with TF-level bootstrap intervals, for attention and for co-expression, expression and degree baselines, also after residualising attention on co-expression, expression and rank distance, and against a degree-preserving permutation null (1,000 rewirings that preserve every TF's out-degree and every target's in-degree). Multiple tests were corrected with the Benjamini--Hochberg procedure~\cite{benjamini1995fdr} or with max-statistic permutation tests. Two decompositions (separating gene-level from pair-specific attention, and comparing attention with co-expression on the same rewired edge sets) were added after the main results had been inspected and are reported as post hoc.

\textbf{Evaluation of the agent runs.} The runs were evaluated from their saved artifacts (prompts, executor scripts, execution outputs and panel critiques) after all runs had finished; no agent code was re-run. Numbers were extracted from the artifacts by script, and the assessment was drafted with LLM assistance and checked by a second agent that re-opened every cited source file and recomputed all counts (S3 File). A panel critique counted as a confirmed defect when the error it described could be confirmed in the cited code lines or outputs, and a defect counted as fixed when the code or outputs of a later execution no longer showed it; corrected code was not re-run independently. Execution failures (crashes and timeouts), which the execution report already showed to the reviewers, were distinguished from defects that only a reading of the code or outputs could reveal. The 305 merged critiques of the 15 executed reviews were not audited exhaustively, so the counts of confirmed defects and of factually wrong critiques are lower bounds, and the panel's precision is not estimated. Each run's final conclusion (that of its last successful iteration where the final execution failed) was compared with the independent analysis on seven claims: attention symmetry, TF$\to$target directionality, recovery of TRRUST edges above chance, recovery above co-expression, robustness of that advantage to degree structure, a pair-specific signal after full adjustment, and causal interpretation. Each claim was marked as agreeing, disagreeing, not addressed, or not directly comparable when the run reached a conclusion about a different estimand. Because the runs chose their own estimands, candidate sets and multiplicity corrections, agreement means the same direction of conclusion, not the same test. Attention symmetry, which was not part of the question, was checked only for whether a run asserted it. Statements that something did not occur (that no run asserted attention symmetry, or that no panel identified a verdict written into the code) rest on regular-expression searches of all executor outputs, generated reports and critiques, supplemented by reading, so a paraphrased mention could have been missed.

\subsection*{Orchestration overhead}

Control-plane overhead was measured through the production runner with a real SQLite database and artifact writing, using the mock backend with zero and 1-s artificial latency and 1--20 concurrent runs (five repeats per level) on an internal and an external drive. Panels in these measurements merged critiques lexically, so the additional adjudication call of the default merge is not included. After every batch, a database integrity check compared iteration counts across database rows, artifact directories and run metadata. Real-backend concurrency was measured in one batch at each of 1, 2, 4 and 8 simultaneous three-lens panels on one benchmark item (gpt-5.6-luna, low effort).

\subsection*{Use of language models outside the evaluated system}

Apart from the models evaluated through the platform (S1 Table), Claude Opus 5.5 (Anthropic; model identifier claude-opus-5-5), used interactively through Claude Code 2.1.284, assisted in writing the benchmark items and their ground truth, the code of the independent case-study analysis, and the synthesis of the case-study runs; separate instances of the same model, working in fresh contexts, checked the benchmark items and the synthesis against their source files. These outputs were validated in different ways. The benchmark was frozen before any reviewer call, and the judges applied its detection criteria; the independent analysis was cross-checked against the agents' own computations of the same quantities (Results); and every number in the synthesis is linked to its source file (S3 File). The same model also assisted with software development of the platform and with drafting and editing this manuscript. Token usage of these sessions is not reported.
